% Importado desde: 2026-03-28_Puente_a_Dirac_3p1_y_Estructura_de_Grupos.tex
\chapter{Derivacion Pedagogica: --- Puente hacia Dirac en \(3+1\) y la estructura de grupos que no debemos perder}
\label{ch:puente-dirac-3p1-grupos}






% verified-by: validators/puentes_grupos.py::test_two_pi_gives_minus_identity
% verified-by: validators/puentes_grupos.py::test_four_pi_gives_identity
% verified-by: validators/puentes_grupos.py::test_rotation_generators_hermitian
% verified-by: validators/puentes_grupos.py::test_boost_generators_antihermitian
% verified-by: validators/puentes_grupos.py::test_six_independent
% verified-by: validators/puentes_grupos.py::test_gamma5_anticommutes
% verified-by: validators/puentes_grupos.py::test_gamma5_squared

\section{Objetivo}

Hasta aqui ya construimos una escalera bastante clara:

\begin{enumerate}[leftmargin=*]
  \item SSH como Dirac en \(1+1\);
  \item grafeno como Dirac en \(2+1\);
  \item semimetales de Dirac/Weyl en \(3D\);
  \item Nielsen--Ninomiya como restriccion estructural;
  \item quantum walks y QCA como ruta hacia una dinamica discreta del espacio-tiempo;
  \item una caminata tipo honeycomb que vuelve a producir Dirac en \(2+1\).
\end{enumerate}

Pero ahora hay que hacer una pausa metodologica importante:

\begin{displayquote}
\textbf{No podemos perder de vista que la meta no es solo \enquote{tener algo tipo Dirac}, sino entender bajo que estructura algebraica y de simetrias puede emerger la ecuacion de Dirac en \(3+1\).}
\end{displayquote}

La pregunta central de este documento es entonces:

\begin{displayquote}
\enquote{Que cambia realmente cuando pasamos de \(2+1\) a \(3+1\), y que estructura de grupos y de algebra debe preservarse o emerger para que el proyecto no se desvie de su objetivo?}
\end{displayquote}

\section{El cambio de perspectiva}

En \(2+1\), muchas veces basta con un Hamiltoniano de dos componentes:
\begin{equation}
H_{2+1}\sim v(\sigma_x p_x + \sigma_y p_y) + m\sigma_z.
\end{equation}

Eso funciona muy bien para grafeno, para quantum walks bidimensionales y para muchas teorias efectivas.

En \(3+1\), en cambio, el problema ya no es solo agregar un tercer momento \(p_z\). Tambien hay que responder:

\begin{enumerate}[leftmargin=*]
  \item que espacio espinorial representa correctamente la teoria;
  \item que matrices implementan la algebra de Clifford adecuada;
  \item que grupo actua sobre esas matrices y sobre los espinores;
  \item como emerge esa estructura a partir de una red o de una dinamica discreta.
\end{enumerate}

\section{Lo que de verdad queremos reproducir}

La ecuacion de Dirac covariante en \(3+1\) es
\begin{equation}
\left(\ii \gamma^\mu \partial_\mu - m\right)\psi = 0,
\label{c17:eq:dirac-cov}
\end{equation}
donde \(\mu=0,1,2,3\) y las matrices \(\gamma^\mu\) satisfacen
\begin{equation}
\{\gamma^\mu,\gamma^\nu\}=2\eta^{\mu\nu}\bbI,
\qquad
\eta^{\mu\nu}=\mathrm{diag}(1,-1,-1,-1).
\label{c17:eq:clifford}
\end{equation}

Si separamos tiempo y espacio, el Hamiltoniano toma la forma
\begin{equation}
\ii\partial_t\psi = H\psi,
\qquad
H = \alpha_x p_x + \alpha_y p_y + \alpha_z p_z + \beta m,
\label{c17:eq:ham-dirac}
\end{equation}
con
\begin{equation}
\alpha_i^2=\beta^2=\bbI,
\qquad
\{\alpha_i,\alpha_j\}=2\delta_{ij}\bbI,
\qquad
\{\alpha_i,\beta\}=0.
\label{c17:eq:alpha-beta}
\end{equation}

Este es el verdadero objetivo algebraico del proyecto.

\section{Primera gran diferencia: dos componentes ya no bastan}

En \(2+1\) pudimos trabajar con matrices de Pauli \(2\times 2\). En \(3+1\), si queremos una teoria de Dirac masiva completa, necesitamos una representacion con \textbf{cuatro componentes}.

La razon es muy concreta:

\begin{enumerate}[leftmargin=*]
  \item necesitamos tres matrices espaciales que anticomuten entre si;
  \item ademas necesitamos una cuarta matriz \(\beta\) que anticomute con las tres;
  \item eso no cabe en una representacion \(2\times 2\).
\end{enumerate}

Una eleccion estandar es
\begin{equation}
\alpha_i = \tau_x \otimes \sigma_i,
\qquad
\beta = \tau_z \otimes \bbI_2,
\label{c17:eq:dirac-rep}
\end{equation}
donde \(\tau_i\) y \(\sigma_i\) son dos copias independientes de matrices de Pauli.

Con esto, el espinor tiene cuatro componentes:
\begin{equation}
\psi=
\begin{pmatrix}
\psi_1\\
\psi_2\\
\psi_3\\
\psi_4
\end{pmatrix}.
\end{equation}

\begin{figure}[ht]
\centering
\begin{tikzpicture}[>=Latex,node distance=1.2cm]
  \node[draw,rounded corners,fill=blue!6,minimum width=2.8cm,minimum height=0.9cm,align=center] (a) {sector \(\tau\)};
  \node[draw,rounded corners,fill=blue!6,minimum width=2.8cm,minimum height=0.9cm,align=center,right=of a] (b) {sector \(\sigma\)};
  \node[draw,rounded corners,fill=orange!10,minimum width=4.8cm,minimum height=1.0cm,align=center,below=1.4cm of $(a)!0.5!(b)$] (c) {espinor de Dirac de \(4\) componentes};
  \draw[->,thick] (a) -- (c);
  \draw[->,thick] (b) -- (c);
\end{tikzpicture}
\caption{Una forma pedagogica de ver el espinor de Dirac en \(3+1\): aparece como combinacion de dos estructuras internas de dos componentes.}
\label{c17:fig:four-components}
\end{figure}

\section{Segunda gran diferencia: aparece una jerarquia de grupos}

Aqui es donde tu observacion sobre la estructura de grupos se vuelve central.

\subsection{Nivel microscopico}

Una red o un automata discreto no tiene de entrada el grupo de Poincare. Lo que suele tener es algo del tipo
\begin{equation}
G_{\mathrm{micro}} = T_{\mathrm{disc}} \rtimes P_{\mathrm{disc}},
\end{equation}
donde:

\begin{enumerate}[leftmargin=*]
  \item \(T_{\mathrm{disc}}\) son traslaciones discretas;
  \item \(P_{\mathrm{disc}}\) es un grupo puntual o un grupo de simetrias locales del paso unitario.
\end{enumerate}

\subsection{Nivel efectivo continuo}

Si el limite infrarrojo es relativista, entonces lo que debe emerger es algo mucho mas grande:
\begin{equation}
\mathcal P = SO(3,1)\ltimes \bbR^{3,1},
\label{c17:eq:poincare}
\end{equation}
el grupo de Poincare.

\subsection{Nivel espinorial}

Pero Dirac no se formula realmente sobre \(SO(3,1)\) a secas, sino sobre su recubrimiento doble
\begin{equation}
SL(2,\bbC).
\label{c17:eq:sl2c}
\end{equation}

Ese es el grupo que actua naturalmente sobre los espinores de Weyl. Y un espinor de Dirac puede verse como
\begin{equation}
\text{Dirac} = \text{Weyl izquierdo} \oplus \text{Weyl derecho}.
\end{equation}

\begin{figure}[ht]
\centering
\begin{tikzpicture}[>=Latex,node distance=1.2cm]
  \node[draw,rounded corners,fill=green!8,minimum width=3.4cm,minimum height=0.9cm,align=center] (m) {simetrias discretas\\de red / QW / QCA};
  \node[draw,rounded corners,fill=green!8,minimum width=3.4cm,minimum height=0.9cm,align=center,below=of m] (l) {simetria efectiva\\de Lorentz};
  \node[draw,rounded corners,fill=green!8,minimum width=3.4cm,minimum height=0.9cm,align=center,below=of l] (s) {\(SL(2,\bbC)\)\\estructura espinorial};
  \node[draw,rounded corners,fill=orange!10,minimum width=3.8cm,minimum height=0.9cm,align=center,right=2.6cm of l] (d) {Dirac en \(3+1\)};
  \draw[->,thick] (m) -- (l);
  \draw[->,thick] (l) -- (s);
  \draw[->,thick] (s) -- (d);
\end{tikzpicture}
\caption{La jerarquia conceptual que no debemos perder. El problema no es solo cinetico: es tambien grupal y espinorial.}
\label{c17:fig:group-ladder}
\end{figure}

\section{Tercera gran diferencia: Weyl y Dirac ya no son lo mismo}

En \(3+1\), la ecuacion de Weyl masiva no existe como teoria independiente simple. Lo natural es:

\begin{enumerate}[leftmargin=*]
  \item un fermion de Weyl sin masa: dos componentes;
  \item dos fermiones de Weyl de quiralidades opuestas;
  \item un termino de masa que los acopla;
  \item el conjunto forma un fermion de Dirac de cuatro componentes.
\end{enumerate}

Este punto es crucial para el proyecto, porque indica un camino muy natural desde construcciones discretas:

\begin{displayquote}
\textbf{Primero construir dinamicas tipo Weyl. Despues entender como acoplar sectores quirales opuestos para formar Dirac en \(3+1\).}
\end{displayquote}

\section{Un modelo minimo para no perder el norte}

Si queremos un objetivo algebraico claro para la siguiente etapa, el Hamiltoniano minimo que debemos tener en la cabeza es
\begin{equation}
H_{\mathrm{target}}(\vq)
=
v_x q_x \alpha_x + v_y q_y \alpha_y + v_z q_z \alpha_z + m\beta,
\label{c17:eq:target}
\end{equation}
con \(\alpha_i,\beta\) satisfaciendo \eqref{c17:eq:alpha-beta}.

En el caso isotropo:
\begin{equation}
H_{\mathrm{target}}(\vq)
=
v\,\vtau_x\otimes(\sigma_x q_x + \sigma_y q_y + \sigma_z q_z)
 + m\,\tau_z\otimes \bbI_2.
\label{c17:eq:target-iso}
\end{equation}

Su espectro es
\begin{equation}
E(\vq)=\pm \sqrt{v^2(q_x^2+q_y^2+q_z^2)+m^2}.
\label{c17:eq:dispersion31}
\end{equation}

Este es el análogo directo del cono de Dirac que vimos antes, pero ahora en tres dimensiones espaciales y con la estructura espinorial correcta.

\begin{figure}[ht]
\centering
\begin{tikzpicture}[>=Latex,node distance=1.25cm]
  \node[draw,rounded corners,fill=blue!6,minimum width=2.5cm,minimum height=0.85cm,align=center] (wx) {\(\alpha_x q_x\)};
  \node[draw,rounded corners,fill=blue!6,minimum width=2.5cm,minimum height=0.85cm,align=center,right=of wx] (wy) {\(\alpha_y q_y\)};
  \node[draw,rounded corners,fill=blue!6,minimum width=2.5cm,minimum height=0.85cm,align=center,right=of wy] (wz) {\(\alpha_z q_z\)};
  \node[draw,rounded corners,fill=orange!10,minimum width=2.5cm,minimum height=0.85cm,align=center,below=1.3cm of wy] (m) {\(\beta m\)};
  \node[draw,rounded corners,fill=green!8,minimum width=4.2cm,minimum height=1.0cm,align=center,below=1.3cm of m] (d) {Dirac \(3+1\)};
  \draw[->,thick] (wx) -- (d);
  \draw[->,thick] (wy) -- (d);
  \draw[->,thick] (wz) -- (d);
  \draw[->,thick] (m) -- (d);
\end{tikzpicture}
\caption{La meta algebraica minima del proyecto. No basta con tres direcciones espaciales: tambien hace falta la matriz de masa y una representacion espinorial adecuada.}
\label{c17:fig:target}
\end{figure}

\section{Que deberia preservar una construccion discreta}

Ahora podemos formular las condiciones que realmente importan.

\subsection{Condicion 1: tres direcciones espaciales independientes}

No podemos quedarnos con una geometria efectivamente bidimensional. La discretizacion debe generar, en el limite efectivo, tres operadores independientes que hagan el papel de \(p_x,p_y,p_z\).

\subsection{Condicion 2: estructura interna suficiente}

La dinamica discreta debe tener suficientes grados de libertad internos para alojar una representacion \(4\times 4\) de la algebra de Clifford \(Cl_{1,3}\).

\subsection{Condicion 3: isotropia efectiva}

Aunque la red microscopica no sea isotropa, la teoria de baja energia debe aproximar
\begin{equation}
v_x \simeq v_y \simeq v_z
\end{equation}
o al menos permitir una reescala controlada para recuperar una forma relativista efectiva.

\subsection{Condicion 4: control de sectores quirales}

Si la construccion parte de bloques tipo Weyl, debe quedar claro:

\begin{enumerate}[leftmargin=*]
  \item cuales son los dos sectores de quiralidad opuesta;
  \item como se acoplan;
  \item bajo que termino aparece la masa de Dirac.
\end{enumerate}

\subsection{Condicion 5: compatibilidad con Nielsen--Ninomiya}

No podemos olvidar el teorema de doblamiento. Cualquier construccion seria en red o en QCA tendra que explicar como aparecen, se emparejan o se controlan los modos extra.

\section{La lectura de grupos en forma compacta}

Podemos resumir el problema de fondo asi:

\begin{center}
\begin{tabular}{@{}lll@{}}
\toprule
Nivel & Objeto & Pregunta clave \\
\midrule
Microscopico & grupo discreto de red/paso & \enquote{Que simetrias exactas tengo?} \\
Efectivo & Lorentz/Poincare & \enquote{Que simetria emerge a baja energia?} \\
Espinorial & \(SL(2,\bbC)\) y Clifford & \enquote{Como se organizan los espinores?} \\
\bottomrule
\end{tabular}
\end{center}

Esta tabla es importante porque evita una confusion comun:

\begin{displayquote}
\enquote{No necesitamos que la red \emph{sea} el grupo de Poincare. Necesitamos que el limite efectivo reconstruya la algebra y las representaciones adecuadas.}
\end{displayquote}

\section{Como se ve el siguiente paso concreto}

Si queremos continuar de forma realmente natural y sin perder la meta, la siguiente etapa deberia tener esta forma:

\begin{enumerate}[leftmargin=*]
  \item construir una \textbf{dinamica discreta tipo Weyl en \(3D\)};
  \item identificar sus simetrias discretas y su espacio interno;
  \item mostrar su limite continuo sin masa;
  \item duplicar o acoplar dos sectores de quiralidad opuesta;
  \item introducir un termino de masa y verificar que aparece el Hamiltoniano de Dirac \eqref{c17:eq:target}.
\end{enumerate}

Esta ruta es mucho mas limpia que intentar saltar directamente desde grafeno a una ecuacion de Dirac completa en \(3+1\).

\section{Mapa conceptual final}

\begin{figure}[ht]
\centering
\begin{tikzpicture}[>=Latex,node distance=1.25cm]
  \node[draw,rounded corners,fill=blue!6,minimum width=3.0cm,minimum height=0.9cm,align=center] (a) {grafeno / QW\\en \(2+1\)};
  \node[draw,rounded corners,fill=blue!6,minimum width=3.0cm,minimum height=0.9cm,align=center,right=of a] (b) {Weyl efectivo\\en \(3D\)};
  \node[draw,rounded corners,fill=green!8,minimum width=3.2cm,minimum height=0.9cm,align=center,right=of b] (c) {acoplamiento\\quiral};
  \node[draw,rounded corners,fill=orange!10,minimum width=3.2cm,minimum height=0.9cm,align=center,right=of c] (d) {Dirac en \(3+1\)};

  \node[draw,rounded corners,fill=purple!8,minimum width=4.3cm,minimum height=0.95cm,align=center,below=1.5cm of b] (e) {grupo de Poincare, \(SL(2,\bbC)\), Clifford};

  \draw[->,thick] (a) -- (b);
  \draw[->,thick] (b) -- (c);
  \draw[->,thick] (c) -- (d);
  \draw[->,thick] (e) -- (b);
  \draw[->,thick] (e) -- (c);
  \draw[->,thick] (e) -- (d);
\end{tikzpicture}
\caption{La estructura del problema a partir de ahora. La teoria de grupos ya no es decoracion: es el esqueleto del paso a \(3+1\).}
\label{c17:fig:final-map}
\end{figure}

\section{Resumen final}

La conclusion de este capitulo es muy simple y muy importante:

\begin{enumerate}[leftmargin=*]
  \item pasar de \(2+1\) a \(3+1\) no es solo agregar una coordenada;
  \item obliga a introducir una estructura espinorial de cuatro componentes;
  \item obliga a reproducir la algebra de Clifford \(Cl_{1,3}\);
  \item obliga a tener presente la jerarquia \(G_{\mathrm{micro}} \to SO(3,1) \to SL(2,\bbC)\);
  \item por eso, la siguiente construccion discreta debe pensarse primero como una teoria tipo Weyl en \(3D\), y solo despues como una teoria de Dirac con masa.
\end{enumerate}

Este es el punto que no podemos perder de vista. A partir de aqui, la investigacion deja de ser una cadena de analogias y pasa a ser un problema de representaciones, simetrias y acoplamientos controlados.
